\documentclass[11pt]{article}

\input{0-parameters.tex}

\begin{document}


\paragraph{Why are nominal prices a problem?} Imagine that interest rates sky rocket due to a crisis, the cost of capital increase, the price of everything increases (without any technical change): is it a productivity loss ? No. However, if we use nominal price in the data we use to calibrate $\zeta$, then it would reflect such a hike in interest prices. 

\paragraph{Outline} We begin by adopting McNerney’s definition of real price growth. We then show how this concept integrates into our theoretical framework. Finally, we assess its empirical implications by examining how it affects the estimated value of $\zeta$.
\tableofcontents

\section{Theory}

 In the McNerney framework, they define the real price increase as  $r_i \equiv \hat{p}_i-\hat{w}_i$ that is ``the rate of change of the real (i.e. inflation-adjusted) price of good $i$.'' They only have one factor here so the entire value added is income. 
We can easily extend this framework, by defining inflation as the average inflation of the factor prices for a specific industry $\hat{r}_i$ (note that we change the notation from McNerney, noting inflation as $\hat{r}_i$ to highlight the fact that it's a growth rate). 


We define the inflation $\hat{r}_i$ as the real input cost inflation faced by sector $i$, weighted by its factor cost shares,
\begin{equation}
\label{eq_def_ri}
\hat{r}_i\equiv \sum_f l_{fi}\hat{w}_f,
\end{equation}
with VA the value added of the sector $i$, $\text{VA}_i=p_iX_i-\sum_j X_{ij}p_j=\sum_f L_{fi}w_f$, $l_{fi}=\frac{L_{fi}w_f}{\text{VA}_i}$ the share of each factor in value added , ensuring that $\sum_f l_{fi}=1$. 


We can re-express the dual definition of TFP growth in real terms, which isolates technical change from inflationary effects (see proof in Section \ref{subsect_derivation_dual_real_price}), 
% \begin{equation}
%     \hat{p}_i-\hat{r}_i=- \hat{A}_i + \sum_j  a_{ji} \left[\hat{p}_j-\hat{r}_i\right].
% \end{equation}

\begin{equation}
\label{eq_dual_real_price}
\hat{A}_i = - \left( (\hat{p}_i-\hat{r}_i) - \sum_{j=1}^N  a_{ji} (\hat{p}_j- \hat{r}_i) \right).
\end{equation}

The new real expression of the dual (Eq. \eqref{eq_dual_real_price}), translates easily into the expression of sector $i$ productivity growth $\hat{A}_i$. Let us define $\hat{\Sigma}_{p_i,\text{real}}\equiv \sum_{j=1}^N  a_{ji} (\hat{p}_j- \hat{r}_i)$. Then, we can then express primal and dual as

\begin{equation}
\hat{A}_i = \hat{X}_i - \hat{\Sigma}_{X_i} = \hat{\Sigma}_{p_i,\text{real}} - (\hat{p}_i-\hat{r}_i).
\label{eq:tfpq_def_SxSp}
\end{equation}
Which means that sales growth is,
\begin{align}
\hat{s}_i &= \hat{X}_i + (\hat{p}_i-\hat{r}_i) = \hat{\Sigma}_{X_i} + \hat{\Sigma}_{p_i,\text{real}}.
\label{eq:siSpSx}
\end{align}
Thus, we define $\zeta$ as 
\begin{equation}
\zeta = \frac{\hat{\Sigma}_{p_i,\text{real}}}{\hat{\Sigma}_{p_i,\text{real}} + \hat{\Sigma}_{X_i}},
\label{eq:zeta}
\end{equation}
which allows us to express sector $i$ productivity growth as a function of output growth and real price growth
\begin{equation}
\label{eq_industry_TFP_real_price}
    \hat{A}_i = \zeta \hat{X}_i + (\zeta - 1) (\hat{p}_i-\hat{r}_i).
\end{equation}


Note that in Eq. \eqref{eq_industry_TFP_real_price}, inflation is industry-specific. It corresponds to the log growth of VA\_PI in KLEMS  (GVA, price indexes (2015)). \ER{However, to simplify the calculations, we chose to use a country-specific inflation}.

Overall, we can then interpret $\hat{p}_i-\hat{r}_i$ as the real price increase of $p_i$ minus the factor inflation and use the exact same framework.


\subsection{Derivation of Real Price Expression}
\label{subsect_derivation_dual_real_price}

In this subsection, we show the derivation of the real price expression of the dual TFP definition  (Eq. \eqref{eq_dual_real_price}).

We start with the definition of real price growth for industry $i$:
\[
\hat{p}_i - \hat{r}_i = \text{nominal price growth} - \text{factor price inflation}
\]

We can express $\hat{p}_i$ using the dual definition of TFP growth, and $\hat{r}_i$ using Eq. \eqref{eq_def_ri}.

\begin{align*}
\hat{p}_i - \hat{r}_i 
&= -\hat{A}_i + \sum_j a_{ji} \hat{p}_j + \sum_f \tilde{l}_{fi} \hat{w}_f - \sum_f \frac{L_{fi} w_f}{\text{VA}_i} \hat{w}_f
\end{align*}
Rearranging terms and using $ \tilde{l}_{fi} = \frac{L_{fi} w_f}{X_i p_i}$ and $\text{VA}_i = p_i X_i - \sum_j X_{ij} p_j$, we get
\begin{align*}
\hat{p}_i - \hat{r}_i &= -\hat{A}_i + \sum_j a_{ji} \hat{p}_j + \sum_f \left( \tilde{l}_{fi} - \frac{L_{fi} w_f}{\text{VA}_i} \right) \hat{w}_f \\
 &= -\hat{A}_i + \sum_j a_{ji} \hat{p}_j + \sum_f \left( \frac{L_{fi} w_f}{X_i p_i} - \frac{L_{fi} w_f}{\text{VA}_i} \right) \hat{w}_f \\
&= -\hat{A}_i + \sum_j a_{ji} \hat{p}_j + \sum_f \left( \frac{1}{X_i p_i} - \frac{1}{p_i X_i - \sum_j X_{ij} p_j} \right) L_{fi} w_f \hat{w}_f \\
&= -\hat{A}_i + \sum_j a_{ji} \hat{p}_j + \sum_f \left( \frac{p_i X_i - \sum_j X_{ij} p_j - X_i p_i}{\text{VA}_i \cdot X_i p_i} \right) L_{fi} w_f \hat{w}_f \\
&= -\hat{A}_i + \sum_j a_{ji} \hat{p}_j + \sum_f \left( \frac{-\sum_j X_{ij} p_j}{\text{VA}_i \cdot X_i p_i} \right) L_{fi} w_f \hat{w}_f \\[1ex]
&= -\hat{A}_i + \sum_j a_{ji} \hat{p}_j + \sum_f \left( \frac{- \frac{\sum_j X_{ij} p_j}{X_i p_i} }{\text{VA}_i} \right) L_{fi} w_f \hat{w}_f \\
\end{align*}

Using that $a_{ji} = \frac{X_{ij} p_j}{X_i p_i} $, we write,
\begin{align*}
\hat{p}_i - \hat{r}_i &= -\hat{A}_i + \sum_j a_{ji} \hat{p}_j + \sum_f \left( \frac{- \sum_j a_{ji} }{\text{VA}_i} \right) L_{fi} w_f \hat{w}_f, \\
\end{align*}
and then rearrange to get an expression akin to the dual definition of TFP growth using real price growth,
\begin{align*}
\hat{p}_i - \hat{r}_i &= -\hat{A}_i + \sum_j a_{ji} \hat{p}_j - \sum_j a_{ji} \sum_f \left( \frac{1}{\text{VA}_i} L_{fi} w_f \hat{w}_f \right) \\
&= -\hat{A}_i + \sum_j a_{ji} \left( \hat{p}_j - \sum_f \frac{L_{fi} w_f}{\text{VA}_i} \hat{w}_f \right) \\
&= -\hat{A}_i + \sum_j a_{ji} \left( \hat{p}_j - \sum_f l_{fi} \hat{w}_f \right) \\
&= -\hat{A}_i + \sum_j a_{ji} \left( \hat{p}_j - \hat{r}_i \right).
\end{align*}

Hence, we have proven Eq. \eqref{eq_dual_real_price}
\begin{equation*}
    \hat{p}_i-\hat{r}_i=- \hat{A}_i + \sum_j  a_{ji} \left[\hat{p}_j-\hat{r}_i\right]
\end{equation*}


\vspace{1em}

\noindent \textbf{Economic interpretation:} Economically, this decomposition tells us that observed price growth minus factor inflation equals the technical change component net of upstream real price effects


\section{Computing $\zeta$ using real prices in KLEMS}

What does it change to use real price growth instead of nominal price growth ? 
Well, quite a lot. 

\paragraph{Distribution of $\zeta$}
We plot the distribution of zeta computed with nominal and real price growth rates.
\begin{figure}[H]
    \centering
     \caption{Distribution of computed zeta in KLEMS}
    \label{g_hist_zeta_real_nominalp}
    \includegraphics[width=0.8\textwidth]{figures/figures_hindcasting/g_hist_zeta_real_nominalp.jpeg}
\end{figure}

\paragraph{Optimisation of $\zeta$ through hindcasting}
To get the optimal value of $\zeta$ by minimising the absolute value of the mean error (AME) between the predicted value and true value of productivity growth at the max horizon $\tau$ (23 years in KLEMS), that is
\begin{equation}
\label{eq_mean_error_t_t+tau}
    \text{AME}_{\tau}=\left|\frac{1}{r}\sum_{i,j,t}(\hat{\tilde{A}}_{i,j,t,\tau}-\hat{A}_{i,j,t,\tau})\right|
\end{equation}
where $r$ is the number of observations (country, sector, pair of years separated by $\tau$ years, with sectoral productivity growth ($\hat{A}_i$) and our framework approximation ($\hat{\tilde{A}}_i$). All countries, industries, years, and horizons (for the growth rates) pooled together.

We then compute the mean error and the root mean square error, at each horizon $\tau$
\[\text{Mean error}_\tau=\frac{1}{n}\sum_i \left(\hat{\tilde{A}}_{t+\tau,i}-\hat{A}_{t+\tau,i}\right)\]
\[\text{RMSE}_\tau=\sqrt{\frac{1}{n}\sum_{i}\left(\hat{\tilde{A}}_{t+\tau,i}-\hat{A}_{t+\tau,i}\right)^2}\]

\paragraph{Results}
With nominal prices, $\zeta^*=0.53$, with real prices $\zeta^*=0.14$.
Errors are pretty similar in both cases, around 200/250 bp at max (figures  \ref{g_mean_error_nominalp} and \ref{g_mean_error_realp}). Using real price growth causes larger RMSE, around 5000bp at the longest horizon (Figure \ref{g_rmse_realp}) compared to nominal price growth (Figure \ref{g_rmse_nominalp}).

There is a significant difference in the values of $\zeta$ using real or nominal prices. Since CCC extra cost are in constant currency and most prices are in real terms, we would consistently \textit{underestimates} the impact of price growth if using a $\zeta$ calibrated on nominal prices. 

\begin{figure}[H]
    \centering
     \caption{Mean and median error on aggregate cumulative productivity - Nominal price growth}
    \label{g_mean_error_nominalp}
    \includegraphics[width=0.8\textwidth]{figures/figures_hindcasting/g_mean_error_nominalp.jpeg}
\end{figure}

\begin{figure}[H]
    \centering
     \caption{Mean and median error on aggregate cumulative productivity - Real price growth}
    \label{g_mean_error_realp}
    \includegraphics[width=0.8\textwidth]{figures/figures_hindcasting/g_mean_error_realp.jpeg}
\end{figure}

\begin{figure}[H]
    \centering
     \caption{RMSE and MAE on aggregate cumulative productivity - Nominal price growth}
    \label{g_rmse_nominalp}
    \includegraphics[width=0.8\textwidth]{figures/figures_hindcasting/g_rmse_nominalp.jpeg}
\end{figure}

\begin{figure}[H]
    \centering
     \caption{RMSE and MAE on aggregate cumulative productivity - Real price growth}
    \label{g_rmse_realp}
    \includegraphics[width=0.8\textwidth]{figures/figures_hindcasting/g_rmse_realp.jpeg}
\end{figure}


\end{document}